\section{Evaluation Scope \& Deployment Limits}
\label{sec:scope-deployment}

Now we will cover the constraints that may affect both the evaluation and deployment of MERIT.

The first constraint lies in the evidence we use, which almost entirely consists of synthetic personas and labels which were engineered by the author
to highlight a subset of specific behaviours for the purpose of a thorough evaluation. In cases such as the adversarial and spearman studies, we attempt
to avoid cherry-picking for success by also cherry-picking for failure. Note that we are not able to use real-world data for the evaluation due to 
the lack of access to real-world data and the need to anonymise the data.

Alongside this, the engines that we use are not commercial products, nor are they open-source alternatives. They are manually created engines,
designed to emulate the scoring logic of commercial products and open-source alternatives under an identical input. Also note that benchmark data
was generated only from the author's workstation, and as such, the exact values are subject to change, based on both the hardware and language runtime.

The final constraint is that MERIT is research-grade. It should never be considered for immediate deployment, with the most significant
blocker being the implementation of LinkedIn ingestion. We must reimplement the ingestion logic to use LinkedIn's official API, rather than the current
prototype which uses Apify's API.

Any further deployment legal, ethical and societal constraints are covered in Section~\ref{sec:ethics-legal-societal}.